\chapter{软件开发基础}

\section{领域导言与演进背景}

\begin{preamblebox}[知识领域概述与变革]
在整个计算机科学的知识版图中，熟练驾驭软件开发过程是每位学子立足的根本。为了最高效地利用计算机解决真实世界问题，学生必须能够流利地阅读、理解并运用多种主流编程语言编写稳健、规范的代码。然而，超越单纯的语法糖和基本编程技能，计算机专业人才更须具备敏锐的数据结构与算法选型能力，并熟练掌握现代软件工程开发与测试工具链。

SDF 知识领域汇聚了软件开发最为核心的基础概念与实操技能，主要面向计算机专业大一（第一学年）基础教学。它涵盖基本编程概念、内置数据结构的工程化应用、高质量编码规范与实践，以及算法对程序运行性能的直观影响。本领域包含的 43 学时核心内容构成了后续进阶专业课的基石。

\textbf{自 CS2013 以来的核心变革：}
\begin{enumerate}
    \item \textbf{强调现代语言原生数据结构赋能}：现代主流语言（如 Python、Java、C++、Rust 等）均内置了极高工业成熟度的关联容器（字典/映射、集合、动态切片）。CS2023 强烈倡导在早期即培养学生直接运用原生/库级高阶数据结构进行工程化解题的能力，避免过早陷入底层指针操作的泥潭。
    \item \textbf{全流程工程规范前置}：将测试驱动意识、单元测试框架、断点调试技术、代码静态审查与可读性重构要求，前置融入大一基础编程阶段。
    \item \textbf{明确划分 SDF 与 AL / SE 的协作边界}：SDF 聚焦于中小型程序的构建、常见数据结构使用与实证性能测试，高阶理论证明（复杂性下界、NPC、形式语义）留待 AL，大规模工业级生命周期管理留待 SE。
\end{enumerate}
\end{preamblebox}

\section{核心学时分配建议}

本领域内容全部为所有计算机专业学生所必需的计算机核心（计算机核心，共 43 学时）：

\begin{table*}[htbp]
\centering\small
\setlength{\tabcolsep}{2.5pt}\footnotesize\begin{tabularx}{0.88\textwidth}{Y c c}
\toprule
\textbf{知识单元} & \textbf{计算机核心} & \textbf{领域核心} \\
\midrule
基础编程概念与实践 & 20 & -- \\
基础数据结构应用 & 12 & -- \\
算法与实证效能 & 6 & -- \\
现代软件开发实践规范 & 5 & -- \\
\midrule
\textbf{合计} & \textbf{43 学时} & \textbf{0 学时} \\
\bottomrule
\end{tabularx}
\end{table*}

\section{知识单元与详细大纲}

\subsection{基础编程概念与实践}
\textbf{【学时建议】}：20 计算机核心学时。

\begin{itemize}
    \item \textbf{程序基础元素}：变量、作用域、生命周期、基本数据类型（整型、浮点型、布尔型、字符型）、表达式求值、赋值语义。
    \item \textbf{程序控制结构}：条件分支（if-else, switch/match 模式匹配）、循环控制、控制流跳转与循环不变量。
    \item \textbf{核心模块化抽象机制}：
    \begin{itemize}
        \item 函数/方法定义：参数传递机制（值传递、引用传递/指针传递、对象引用共享）、默认参数与可变参数。
        \item 作用域与命名空间：局部变量、全局变量、闭包上下文与变量遮蔽。
        \item 面向对象基础：类与对象的定义、数据封装、访问控制修饰符、简单继承与多态。
    \end{itemize}
    \item \textbf{现代复合数据类型}：序列结构（静态数组、动态数组/列表）、键值映射（字典/Map/哈希表）、集合（Set）、元组、不可变数据结构。
    \item \textbf{标准输入输出与外部交互}：控制台标准流 I/O、磁盘持久化文件读写、RESTful Web API 资源获取与 JSON 结构化解析。
    \item \textbf{递归思想与实现}：递归基准条件、递归工作栈、树状递归与递归向迭代转换。
    \item \textbf{运行时异常防御与处理}：异常抛出、捕获、异常层级与资源优雅释放。
    \item \textbf{代码心智模型与执行追踪}：堆栈内存布局、对象引用图谱、代码执行控制流跟踪。
\end{itemize}

\subsection{基础数据结构应用}
\textbf{【学时建议】}：12 计算机核心学时。

\begin{itemize}
    \item \textbf{标准抽象数据类型 (ADT) 运用}：列表（List）、栈（Stack，LIFO）、队列（Queue，FIFO）、双端队列、集合（Set）与字典（Dictionary）。
    \item \textbf{文本与字符串高级处理}：字符串不可变性、子串查找、正则表达式匹配与文本流切分。
    \item \textbf{针对业务场景的数据结构选型}：分析特定场景（高频写入 vs 频繁随机访问）下各种内置容器的渐进开销与内存布局。
    \item \textbf{实证性能度量}：编写基准评测代码，度量不同数据规模下容器操作的物理耗时。
\end{itemize}

\subsection{算法与实证效能}
\textbf{【学时建议】}：6 计算机核心学时。

\begin{itemize}
    \item \textbf{算法本质与效能度量}：算法定义、正确性验证直觉、时间与空间复杂性直观感知。
    \item \textbf{入门经典算法族}：
    \begin{itemize}
        \item 常用排序：选择排序、插入排序、快速排序思想与库级排序函数的使用。
        \item 经典查找：线性搜索、有序序列二分查找及其死循环规避。
        \item 树与图基础遍历：深度优先遍历（DFS）与广度优先遍历（BFS）的直观递归与队列实现。
    \end{itemize}
    \item \textbf{实测性能对比}：在真实硬件上对 $O(n)$ 与 $O(\log n)$、$O(n^2)$ 与 $O(n \log n)$ 算法进行规模压力评测。
\end{itemize}

\subsection{现代软件开发实践规范}
\textbf{【学时建议】}：5 计算机核心学时。

\begin{itemize}
    \item \textbf{工业级编码风格与可读性}：Clean Code 规范、自解释的变量与函数命名、模块职责单一原则、适度而精准的内联注释。
    \item \textbf{现代集成开发环境 (IDE) 与调试技术}：代码导航、符号查找、断点设置、条件断点、单步执行、变量监视与调用栈检查。
    \item \textbf{单元测试与测试用例设计}：基于现代测试框架（JUnit、pytest 等）编写断言；等价类划分、边界值分析与异常分支测试用例构造。
    \item \textbf{测试局限性反思}：深刻理解“测试可以证明程序存在缺陷，但永远无法证明程序没有缺陷”（Dijkstra 箴言）。
\end{itemize}

\section{专业素养与思维倾向}

\begin{itemize}
    \item \textbf{自主解决问题的韧性}：遇到语法或逻辑报错时，善于自主查阅官方文档、API 参考与技术社区排查方案。
    \item \textbf{实验探索精神}：勇于编写小 Demo 验证未知语言特性，熟练运用调试器单步推演探索程序行为。
    \item \textbf{对技术底层机理的好奇心}：不仅满足于代码能跑通，更渴望理解 IDE 内部运行机制、内存分配与编译流程。
    \item \textbf{严密的系统性与细密性 (Systematic \& Meticulous)}：在处理边界空指针、数组下标越界与输入非法校验时保持严谨。
\end{itemize}

\section{课程设计与教学实施建议}

联合工作组建议将 SDF 模块有机分布在高校第一学年的经典两门课中：
\begin{enumerate}
    \item \textbf{CS1：程序设计基础 I (Introductory Programming)}：
    覆盖 20 学时编程概念、12 学时数据结构应用、2 学时算法入门、2 学时测试调试实践及 2--4 学时工程伦理与规范。目标是使学生能独立开发包含函数抽象、对象交互与文件/网络 API 调用的百行级稳健程序。
    \item \textbf{CS2：数据结构与面向对象实战 (Data Structures \& OOP)}：
    深化承接 SDF 的后续 4 学时高级算法（与 AL 衔接）和 3 学时高阶软件工程规范（与 SE 衔接）。
\end{enumerate}

\section{真实能力任务规范}

\begin{taskbox}[任务 1：模块化功能组件的规范开发与单元测试]
\textbf{【任务情境】}：根据给定的业务功能规约（例如多格式日志解析器），设计并编写一个独立的函数或类。\\
\textbf{【能力陈述】}：编写高可读性、高聚合度的代码，覆盖所有入参边界，并编写单元测试验证所有分支覆盖。\\
\textbf{【能力领域】}：软件开发 \quad \textbf{【能力单元】}：系统实现与质量测试\\
\textbf{【达标等级】}：\textbf{开发}
\end{taskbox}

\begin{taskbox}[任务 2：复杂逻辑缺陷的断点调试与定位修复]
\textbf{【任务情境】}：某学生在线考试评分系统中存在偶发性并发计分错误和异常越界崩溃。\\
\textbf{【能力陈述】}：运用现代 IDE 的断点、调用栈溯源与变量跟踪功能，在没有源码注释的情况下快速定位 Bug 根源并重构修复。\\
\textbf{【能力领域】}：软件开发 \quad \textbf{【能力单元】}：系统调试与维护\\
\textbf{【达标等级】}：\textbf{评估与开发}
\end{taskbox}

\begin{taskbox}[任务 3：同侪代码评审 与代码异味排查]
\textbf{【任务情境】}：在敏捷团队中对同组开发者提交的代码拉取请求（PR）进行严格的同行评审。\\
\textbf{【能力陈述】}：阅读并理解第三方代码，准确指出其中的潜在逻辑漏洞、死循环边界隐患、不规范命名与可维护性缺陷，并给出具体的重构建议。\\
\textbf{【能力领域】}：软件工程与应用 \quad \textbf{【能力单元】}：代码评审与改进\\
\textbf{【达标等级】}：\textbf{评估}
\end{taskbox}

\begin{taskbox}[任务 4：基于第三方 RESTful API 的轻量级数据聚合程序]
\textbf{【任务情境】}：构建一个实时天气或交通状况聚合展示终端应用。\\
\textbf{【能力陈述】}：熟练调用网络请求库发起 HTTP 访问，安全处理网络超时、鉴权失败及 JSON 反序列化异常，将数据缓存至本地哈希表中并做格式化呈现。\\
\textbf{【能力领域】}：软件与应用 \quad \textbf{【能力单元】}：API 集成开发\\
\textbf{【达标等级】}：\textbf{开发}
\end{taskbox}

\section{附录：核心主题与技能要求等级对照表}

\begin{longtable}{p{0.20\textwidth} p{0.48\textwidth} c c c}
\toprule
\textbf{知识单元} & \textbf{核心知识点与具体主题} & \textbf{认知等级} & \textbf{分类} & \textbf{学时} \\
\midrule
\endfirsthead
\toprule
\textbf{知识单元} & \textbf{核心知识点与具体主题} & \textbf{认知等级} & \textbf{分类} & \textbf{学时} \\
\midrule
\endhead
\bottomrule
\endfoot
\multirow{2}{*}{\makecell[l]{软件开发基础/基础编程概念\\与实践}} 
& 变量、表达式、分支迭代控制、函数封装、异常与 I/O & 开发 & 计算机核心 & 18 \\
& 程序错误类型、断点调试技术、代码跟踪与阅读理解 & 应用 & 计算机核心 & 2 \\
\midrule
\multirow{2}{*}{\makecell[l]{软件开发基础/基础数据结构}} 
& 常用容器（列表、栈、队列、字典、集合）应用与字符串处理 & 开发 & 计算机核心 & 10 \\
& 容器选型权衡、时空性能实测与内存开销评估 & 评估 & 计算机核心 & 2 \\
\midrule
\multirow{2}{*}{\makecell[l]{软件开发基础/算法}} 
& 算法概念、基础查找与排序、树图简单遍历 & 解释 & 计算机核心 & 4 \\
& 算法时空效率的物理测量与规模敏感性实证评测 & 评估 & 计算机核心 & 2 \\
\midrule
\multirow{3}{*}{\makecell[l]{软件开发基础/软件开发实践}} 
& 工业级编码规范、代码风格与可读性重构 & 评估 & 计算机核心 & 1 \\
& 单元测试框架使用与测试用例设计（等价类与边界值） & 开发 & 计算机核心 & 2 \\
& 现代 IDE 工具链（断点调试、性能分析）熟练使用 & 应用 & 计算机核心 & 2 \\
\end{longtable}
